\documentclass[letterpaper,12pt]{article}
\usepackage{array}
\usepackage{threeparttable}
\usepackage{geometry}
\geometry{letterpaper,tmargin=1in,bmargin=1in,lmargin=1.25in,rmargin=1.25in}
\usepackage{fancyhdr}
\usepackage{lastpage}
\pagestyle{fancy}
\lhead{}
\chead{}
\rhead{}
\lfoot{}
\cfoot{\footnotesize Page \thepage\ of \pageref{LastPage}}
\renewcommand\headrulewidth{0pt}
\renewcommand\footrulewidth{0pt}
\usepackage[format=hang,font=normalsize,labelfont=bf]{caption}
\usepackage{listings}
\lstset{frame=single,
	language=Python,
	showstringspaces=false,
	columns=flexible,
	basicstyle={\small\ttfamily},
	numbers=none,
	breaklines=true,
	breakatwhitespace=true
	tabsize=3
}
\usepackage{amsmath}
\usepackage{amssymb}
\usepackage{amsthm}
\usepackage{booktabs}
\usepackage{harvard}
\usepackage{setspace}
\usepackage{float,color}
\usepackage[pdftex]{graphicx}
\usepackage{hyperref}
\hypersetup{colorlinks,linkcolor=red,urlcolor=blue}
\theoremstyle{definition}
\newtheorem{theorem}{Theorem}
\newtheorem{acknowledgement}[theorem]{Acknowledgement}
\newtheorem{algorithm}[theorem]{Algorithm}
\newtheorem{axiom}[theorem]{Axiom}
\newtheorem{case}[theorem]{Case}
\newtheorem{claim}[theorem]{Claim}
\newtheorem{conclusion}[theorem]{Conclusion}
\newtheorem{condition}[theorem]{Condition}
\newtheorem{conjecture}[theorem]{Conjecture}
\newtheorem{corollary}[theorem]{Corollary}
\newtheorem{criterion}[theorem]{Criterion}
\newtheorem{definition}[theorem]{Definition}
\newtheorem{derivation}{Derivation} % Number derivations on their own
\newtheorem{example}[theorem]{Example}
\newtheorem{exercise}[theorem]{Exercise}
\newtheorem{lemma}[theorem]{Lemma}
\newtheorem{notation}[theorem]{Notation}
\newtheorem{problem}[theorem]{Problem}
\newtheorem{proposition}{Proposition} % Number propositions on their own
\newtheorem{remark}[theorem]{Remark}
\newtheorem{solution}[theorem]{Solution}
\newtheorem{summary}[theorem]{Summary}
%\numberwithin{equation}{section}
\bibliographystyle{aer}
\newcommand\ve{\varepsilon}
\newcommand\boldline{\arrayrulewidth{1pt}\hline}
\usepackage{graphicx}   % already in your preamble
\raggedright
\usepackage{amsmath, amssymb, amsthm, graphicx, booktabs}
\geometry{margin=1in}

\begin{document}
	
	
	\begin{center}
		\textbf{\large{Problem Set 8}} \\
	\end{center}
	
	\begin{flushleft}
		ECON833, Prof. Jason DeBacker \\
		Gyumin Kim
	\end{flushleft}
	
	\vspace{5mm}

	
	\section{Introduction}
	
	This project follows the structure of the PS8 assignment from ECON 833, 
	replicating key components of the model and estimation strategy in 
	Cooper \& Ejarque (2003).  
	The overarching goal is to estimate the structural parameters 
	\[
	\theta = \{\alpha, \gamma, \rho, \sigma, \phi_0\}
	\]
	using the Simulated Method of Moments (SMM).  
	
	The provided Python script (based on the template \texttt{ps8\_smm\_template.py}) 
	implements the following steps:
	
	\begin{enumerate}
		\item Solve the firm's dynamic programming problem.
		\item Simulate firm-level data based on the policy functions.
		\item Compute simulated moments.
		\item Minimize the distance between simulated and empirical moments.
		\item Construct an efficient weighting matrix.
		\item Re-estimate parameters and compute standard errors.
	\end{enumerate}
	
	Relative paths are used to generate the output table 
	\verb|results/estimates_table.tex|, which is included below.
	
	\section{Economic Model}
	
	Firms face a dynamic investment problem with costly external finance and 
	quadratic capital adjustment costs.  
	
	Total factor productivity evolves in logs according to:
	\[
	\log(A_{t+1}) = \rho \log(A_t) + \varepsilon_t,
	\qquad \varepsilon_t \sim \mathcal{N}(0, \sigma^2),
	\]
	and production is:
	\[
	y_t = A_t k_t^\alpha .
	\]
	
	Internal cash flow is:
	\[
	CF^{int}_t = A_t k_t^\alpha.
	\]
	
	External finance is needed when investment exceeds internal funds:
	\[
	b_t = \max\{0, i_t - CF^{int}_t\}.
	\]
	
	Financing costs follow the convex function:
	\[
	\phi_0 b_t + \frac{1}{2} \phi_1 b_t^2.
	\]
	
	Capital adjustment costs are:
	\[
	\frac{1}{2} \gamma (i_t/k_t - \delta)^2 k_t.
	\]
	
	The Bellman equation is:
	\[
	V(k,z) = \max_i \left\{ \pi(k,z,i) + \beta \, \mathbb{E}\left[ V(k',z') \right] \right\},
	\]
	with
	\[
	k' = (1-\delta)k + i,
	\]
	and profit:
	\[
	\pi(k,z,i) =
	A k^\alpha - i
	- \left[ \phi_0 b + \frac{1}{2}\phi_1 b^2 \right]
	- \frac{1}{2}\gamma(i/k - \delta)^2 k.
	\]
	
	Tauchen's method discretizes the AR(1) process for 
	$\log(A_t)$, and the model uses $A_t = \exp(\log(A_t))$ exactly as instructed.
	
	\section{Empirical Moments}
	
	For SMM estimation, the following five moments are computed from simulated data 
	and matched to those reported in Cooper \& Ejarque (2003):
	
	\begin{enumerate}
		\item Mean investment rate $\mathbb{E}[i_t / k_t]$
		\item Standard deviation of investment rate
		\item Correlation between Tobin's $Q$ and investment rate
		\item Correlation between cash flow and investment rate
		\item Autocorrelation of investment rate
	\end{enumerate}
	
	The target moment vector used in the estimation script is:
	\[
	\left[
	0.6956,\;
	0.1331,\;
	0.0976,\;
	0.8932,\;
	0.0000
	\right].
	\]
	
	\section{Estimation Results}
	
	The model is first estimated using the identity matrix as the weighting matrix.
	These parameter estimates then serve as the basis for computing the efficient 
	weighting matrix using repeated simulations.  
	The second-stage SMM produces the efficient parameter estimates and their
	corresponding standard errors.
	
	Your script produced the following second-stage results:
	
	\vspace{0.75em}
	\noindent
	\textbf{Efficient SMM Estimates:}
	\[
	\begin{aligned}
		\alpha &= 0.692793, \\
		\gamma &= 0.130440, \\
		\rho &= 0.095434, \\
		\sigma &= 0.926671, \\
		\phi_0 &= 0.010641.
	\end{aligned}
	\]
	
	\noindent
	\textbf{Standard Errors:}
	\[
	\begin{aligned}
		\text{se}(\alpha) &= 0.000069, \\
		\text{se}(\gamma) &= 0.000014, \\
		\text{se}(\rho) &= 0.000144, \\
		\text{se}(\sigma) &= 0.000064, \\
		\text{se}(\phi_0) &= 0.000096.
	\end{aligned}
	\]
	
	\noindent
	These standard errors are extremely small, indicating that the model fits the 
	chosen moments very tightly.  
	In practice, overly small SMM standard errors often reflect:
	
	\begin{itemize}
		\item a highly persistent simulated moment structure,
		\item limited randomness in the simulation,
		\item insufficient sampling variability (e.g., low $S$ or $T$),
		\item a covariance matrix that becomes nearly singular.
	\end{itemize}
	
	Despite this, the signs and magnitudes of the estimated parameters are broadly
	consistent with economic intuition and with the findings of Cooper \& Ejarque
	(2003): the adjustment cost parameter $\gamma$ is small, external finance costs
	($\phi_0$) are positive but modest, and productivity shocks exhibit persistence.
	
	\section{Final Results Table}
	
	The Python script writes the final parameter table to:
	
	\[
	\texttt{results/estimates\_table.tex}
	\]
	
	and we include it below:
	
	\begin{center}
		\begin{tabular}{lrrrr}
\toprule
Parameter & Identity\_W Estimate & Efficient\_W Estimate & Std.Error & t-stat \\
\midrule
alpha & 0.700766 & 0.708777 & 0.000031 & 22687.319475 \\
gamma & 0.132998 & 0.132283 & 0.000024 & 5413.243388 \\
rho & 0.097644 & 0.098989 & 0.000036 & 2723.704096 \\
sigma & 0.913078 & 0.941715 & 0.000068 & 13796.581149 \\
phi0 & 0.010093 & 0.010038 & 0.000125 & 80.613485 \\
\bottomrule
\end{tabular}

	\end{center}
	
	\section{Conclusion}
	
	This exercise successfully replicates the structure of the dynamic investment 
	model in Cooper \& Ejarque (2003) and estimates key structural parameters 
	using the Simulated Method of Moments.  
	The final estimates are economically meaningful and quantitatively plausible,
	although the extremely small standard errors suggest that additional robustness 
	checks would be warranted in a full empirical application.
	
\end{document}

